\chapter{مبانی و فناوری‌های مورد استفاده}
\thispagestyle{empty}

\section{مقدمه}
توسعه یک سامانه پرسش و پاسخ پیشرفته و ایمن پزشکی، نیازمند بهره‌گیری از ابزارها و فناوری‌های متنوعی در حوزه‌های پردازش زبان، پایگاه‌های داده و مهندسی نرم‌افزار است. در این فصل، سه فناوری و چارچوب کلیدی مورد استفاده در پیاده‌سازی این سامانه یعنی لنگ‌چین، لنگ‌گراف و کیودرانت معرفی و بررسی می‌شوند.

\section{لنگ‌چین \lr{(LangChain)}}
لنگ‌چین یک چارچوب متن‌باز و قدرتمند است که به منظور تسهیل توسعه برنامه‌های کاربردی مبتنی بر مدل‌های زبانی بزرگ طراحی شده است. این چارچوب ابزارها و رابط‌های استانداردی را برای اتصال مدل‌های زبانی به منابع داده خارجی، ابزارهای جستجو، و حافظه فراهم می‌کند. در توسعه سامانه‌های پرسش و پاسخ مبتنی بر \lr{RAG}، لنگ‌چین نقش زیرساخت اصلی را برای مدیریت زنجیره‌های پردازشی \lr{(Chains)}، قالب‌بندی پرامپت‌ها و هماهنگ‌سازی اجزای مختلف سیستم ایفا می‌کند.

\section{لنگ‌گراف \lr{(LangGraph)}}
به منظور مدیریت پیچیدگی‌های مربوط به درک پرسش، ترجمه، جستجو و تولید پاسخ، از رویکرد سیستم‌های چندعاملی \lr{(Multi-Agent)} استفاده شده است. چارچوب لنگ‌گراف که به عنوان توسعه‌ای بر روی لنگ‌چین معرفی شده است، ابزاری قدرتمند برای طراحی گردش‌کارهای پیچیده و مبتنی بر وضعیت \lr{(Stateful)} است. این ابزار امکان تعریف گراف‌های چرخه‌ای را فراهم می‌کند که در آن عامل‌های مختلف (مانند عامل مترجم، بازیاب و مسیر‌یاب) می‌توانند به صورت هماهنگ و در چرخه‌های تکرارپذیر با یکدیگر تعامل داشته باشند تا پاسخ نهایی با بالاترین کیفیت آماده شود.

\section{پایگاه داده برداری کیودرانت \lr{(Qdrant)}}
در معماری \lr{RAG}، متون و اسناد پزشکی پس از پردازش، باید به شکل بردارهای عددی (جاسازی یا \lr{Embeddings}) ذخیره شوند تا امکان جستجوی معنایی فراهم گردد. کیودرانت یک موتور جستجوی برداری متن‌باز و بسیار کارآمد است که منحصراً برای ذخیره‌سازی و جستجوی سریع در فضای ابعادی بالا طراحی شده است. استفاده از پایگاه داده \lr{Qdrant} به سامانه اجازه می‌دهد تا در میان حجم انبوهی از متون تخصصی اورولوژی، مرتبط‌ترین قطعات متنی با پرسش کاربر را با سرعت و دقت بسیار بالا بازیابی کند.